% GENERATED FILE. Edit canonical lessons, then run npm run generate:books.
% canonical-source-hash: fnv1a64:d15096988fd590eb
% canonical-lessons: MW-C275-saravno, MW-C275-sabas, MW-C275-inam, MW-C275-jas, MW-C275-ray

\chapter{Well Done!}
\label{ch:mw-well-done}
\hlchaptermodality{275}

\begin{chapteropening}
\textbf{By the end of this chapter:} \emph{I can praise someone and ask what they think.}

Complete the last lesson of chapter 275: I can praise someone and ask what they think.
\end{chapteropening}

\section[\texorpdfstring{\mw{सरावणो} (sarāvṇo)}{सरावणो (sarāvṇo)}]{\mw{सरावणो} (sarāvṇo) — to praise}

\label{lesson:MW-C275-saravno}



\begin{quote}
Before the new one: say the Marwadi for mercy, then the Marwadi for alms.
\end{quote}

\subsection*{You'll want to know: \mw{सरावणो}}
\textbf{\mw{सरावणो}} — \emph{sarāvṇo} — \textquotedblleft{}to praise\textquotedblright{}.

\begin{quote}
\mw{सगळा} \mw{गीत} \mw{सरावै} \mw{है।}
\end{quote}

\begin{quote}
\emph{sagḷā gīt sarāvai hai.}
\end{quote}

\begin{quote}
Everyone praises the song.
\end{quote}

Past (I …): \mw{म्हैं} \mw{गीत} \mw{सरायो}, \emph{mhaĩ gīt sarāyo}.

\begin{cousinweb}
\textbf{\mw{सरावणो}} (\emph{sarāvṇo}) is inherited from Sanskrit \emph{ślāghate}, \textquotedblleft{}praises\textquotedblright{}; CDIAL lists Old Marwari \emph{sarāi}, \textquotedblleft{}praises\textquotedblright{} (CDIAL 12734).
\end{cousinweb}

\begin{grammarlens}[title={What you've built}]
One more word for asking and answering.
\end{grammarlens}

\subsection*{Your turn}
\begin{itemize}
\raggedright
  \item \emph{Say it:} \emph{sarāvṇo}
  \item \emph{Say it:} \emph{sarāvṇo}, once more
  \item \emph{Say it:} \emph{sagḷā gīt sarāvai hai.} — Everyone praises the song.
  \item \emph{Recall:} say the Marwadi for mercy, then the Marwadi for alms, then say \emph{sarāvṇo} again
\end{itemize}

\begin{tcolorbox}[breakable,colback=teal!4,colframe=teal!35!black,title={Before you move on}]
What does \mw{सरावणो} mean? (\textquotedblleft{}To praise\textquotedblright{}.) Say it once more.
\end{tcolorbox}
\section[\texorpdfstring{\mw{साबास} (sābās)}{साबास (sābās)}]{\mw{साबास} (sābās) — well done}

\label{lesson:MW-C275-sabas}



\begin{quote}
Before the new one: say the Marwadi for alms, then the Marwadi for to praise.
\end{quote}

\subsection*{You'll want to know: \mw{साबास}}
\textbf{\mw{साबास}} — \emph{sābās} — \textquotedblleft{}well done\textquotedblright{}.

\begin{quote}
\mw{साबास}, \mw{छोरी}!
\end{quote}

\begin{quote}
\emph{sābās, chhorī!}
\end{quote}

\begin{quote}
Well done, girl!
\end{quote}

\begin{cousinweb}
\textbf{\mw{साबास}} (\emph{sābās}) is usually traced to Persian \emph{shābāsh}, \textquotedblleft{}bravo\textquotedblright{}; it is said on its own, as praise.
\end{cousinweb}

\begin{grammarlens}[title={What you've built}]
One more word for asking and answering.
\end{grammarlens}

\subsection*{Your turn}
\begin{itemize}
\raggedright
  \item \emph{Say it:} \emph{sābās}
  \item \emph{Say it:} \emph{sābās}, once more
  \item \emph{Say it:} \emph{sābās, chhorī!} — Well done, girl!
  \item \emph{Recall:} say the Marwadi for alms, then the Marwadi for to praise, then say \emph{sābās} again
\end{itemize}

\begin{tcolorbox}[breakable,colback=teal!4,colframe=teal!35!black,title={Before you move on}]
What does \mw{साबास} mean? (\textquotedblleft{}Well done\textquotedblright{}.) Say it once more.
\end{tcolorbox}
\section[\texorpdfstring{\mw{इनाम} (inām)}{इनाम (inām)}]{\mw{इनाम} (inām) — a prize}

\label{lesson:MW-C275-inam}



\begin{quote}
Before the new one: say the Marwadi for to praise, then the Marwadi for well done.
\end{quote}

\subsection*{You'll want to know: \mw{इनाम}}
\textbf{\mw{इनाम}} — \emph{inām} — \textquotedblleft{}a prize\textquotedblright{}.

\begin{quote}
\mw{छोरी} \mw{इनाम} \mw{जीतै} \mw{है।}
\end{quote}

\begin{quote}
\emph{chhorī inām jītai hai.}
\end{quote}

\begin{quote}
The girl wins a prize.
\end{quote}

A masculine noun.

\begin{cousinweb}
\textbf{\mw{इनाम}} (\emph{inām}) is usually traced, through Persian \emph{inʿām}, to Arabic \emph{inʿām}, \textquotedblleft{}a gift, a favour\textquotedblright{}.
\end{cousinweb}

\begin{grammarlens}[title={What you've built}]
One more word for asking and answering.
\end{grammarlens}

\subsection*{Your turn}
\begin{itemize}
\raggedright
  \item \emph{Say it:} \emph{inām}
  \item \emph{Say it:} \emph{inām}, once more
  \item \emph{Say it:} \emph{chhorī inām jītai hai.} — The girl wins a prize.
  \item \emph{Recall:} say the Marwadi for to praise, then the Marwadi for well done, then say \emph{inām} again
\end{itemize}

\begin{tcolorbox}[breakable,colback=teal!4,colframe=teal!35!black,title={Before you move on}]
What does \mw{इनाम} mean? (\textquotedblleft{}A prize\textquotedblright{}.) Say it once more.
\end{tcolorbox}
\section[\texorpdfstring{\mw{जस} (jas)}{जस (jas)}]{\mw{जस} (jas) — credit; renown}

\label{lesson:MW-C275-jas}



\begin{quote}
Before the new one: say the Marwadi for well done, then the Marwadi for a prize.
\end{quote}

\subsection*{You'll want to know: \mw{जस}}
\textbf{\mw{जस}} — \emph{jas} — \textquotedblleft{}credit; renown\textquotedblright{}.

\begin{quote}
\mw{यो} \mw{जस} \mw{थारो} \mw{है।}
\end{quote}

\begin{quote}
\emph{yo jas thāro hai.}
\end{quote}

\begin{quote}
The credit is yours.
\end{quote}

A masculine noun.

\begin{cousinweb}
\textbf{\mw{जस}} (\emph{jas}) is inherited from Sanskrit \emph{yaśas}, \textquotedblleft{}splendour, renown\textquotedblright{}; CDIAL gives no Marwari form in this entry, but lists Hindi \emph{jas}, \textquotedblleft{}fame\textquotedblright{} (CDIAL 10443).
\end{cousinweb}

\begin{grammarlens}[title={What you've built}]
One more word for asking and answering.
\end{grammarlens}

\subsection*{Your turn}
\begin{itemize}
\raggedright
  \item \emph{Say it:} \emph{jas}
  \item \emph{Say it:} \emph{jas}, once more
  \item \emph{Say it:} \emph{yo jas thāro hai.} — The credit is yours.
  \item \emph{Recall:} say the Marwadi for well done, then the Marwadi for a prize, then say \emph{jas} again
\end{itemize}

\begin{tcolorbox}[breakable,colback=teal!4,colframe=teal!35!black,title={Before you move on}]
What does \mw{जस} mean? (\textquotedblleft{}Credit; renown\textquotedblright{}.) Say it once more.
\end{tcolorbox}
\section[\texorpdfstring{\mw{राय} (rāy)}{राय (rāy)}]{\mw{राय} (rāy) — an opinion}

\label{lesson:MW-C275-ray}



\begin{quote}
Before the new one: say the Marwadi for a prize, then the Marwadi for credit for a deed.
\end{quote}

\subsection*{You'll want to know: \mw{राय}}
\textbf{\mw{राय}} — \emph{rāy} — \textquotedblleft{}an opinion\textquotedblright{}.

\begin{quote}
\mw{थारी} \mw{राय} \mw{कांई} \mw{है}?
\end{quote}

\begin{quote}
\emph{thārī rāy kāĩ hai?}
\end{quote}

\begin{quote}
What's your opinion?
\end{quote}

A feminine noun.

\begin{cousinweb}
\textbf{\mw{राय}} (\emph{rāy}) is usually traced, through Persian \emph{rāy}, to Arabic \emph{raʾy}, \textquotedblleft{}opinion, view\textquotedblright{}. It is feminine: \textbf{\mw{थारी} \mw{राय}}, your opinion.
\end{cousinweb}

\begin{grammarlens}[title={What you've built}]
One more word for asking and answering.
\end{grammarlens}

\subsection*{Your turn}
\begin{itemize}
\raggedright
  \item \emph{Say it:} \emph{rāy}
  \item \emph{Say it:} \emph{rāy}, once more
  \item \emph{Say it:} \emph{thārī rāy kāĩ hai?} — What's your opinion?
  \item \emph{Recall:} say the Marwadi for a prize, then the Marwadi for credit for a deed, then say \emph{rāy} again
\end{itemize}

\begin{tcolorbox}[breakable,colback=teal!4,colframe=teal!35!black,title={Before you move on}]
What does \mw{राय} mean? (\textquotedblleft{}An opinion\textquotedblright{}.) Say it once more.
\end{tcolorbox}
